\documentclass[11pt]{article}
\usepackage[margin=1in]{geometry}
\usepackage{graphicx}
\usepackage{amsmath}
\usepackage{amssymb}
\usepackage{booktabs}
\usepackage{natbib}
\usepackage{hyperref}
\usepackage{xcolor}
\usepackage{url}
\bibliographystyle{plainnat}

\title{Cyclic Homogeneous Oscillation Detection:\\ An Auto-correlation-Based Method for Identifying Non-sinusoidal Neural Oscillations in ECoG, EEG, and SEEG}
\author{Anonymous}
\date{}

\begin{document}
\maketitle

\begin{abstract}
Detecting when and where neural oscillations occur, and what frequency they truly occupy, is a prerequisite for almost every study of dynamic brain function. The standard approach---identifying peaks over a $1/f$ aperiodic background in a power spectrum---assumes that oscillations are sinusoidal and stationary. This assumption is violated by mu rhythms, hippocampal theta, and cortical beta, all of which are non-sinusoidal and therefore generate harmonic peaks in the frequency domain that are indistinguishable, to an FFT, from independent rhythms. We define three criteria that a neural oscillation must satisfy---a peak over $1/f$ noise in the time and frequency domains, at least two full cycles, and agreement between the spectral center frequency and the periodicity of the signal's auto-correlation---and present Cyclic Homogeneous Oscillation (CHO) detection, which implements these criteria in a single pipeline. We validate CHO on synthetic non-sinusoidal bursts of 2.5 cycles (1--3 s) convolved with $1/f$ noise, and on four human datasets: ECoG recordings from 8 subjects and scalp EEG from 7 subjects during the pre-stimulus period of an auditory reaction-time task, resting-state ECoG from 6 subjects, and resting-state stereo-EEG from 6 subjects with hippocampal depth electrodes. At the signal-to-noise ratios observed in our physiological data ($-7$~dB for EEG, $-6$~dB for ECoG), CHO detects 50--60\% of true oscillations and substantially outperforms FOOOF, OEvent, and SPRiNT in specificity and peak-frequency accuracy. In resting-state recordings, CHO identifies 11~Hz (450~ms) oscillations in primary auditory cortex, 7~Hz (450~ms) with prominent $>$500~ms beta in primary motor cortex, 17~Hz beta in Broca's area, and 8~Hz (450~ms) in hippocampus, with hippocampal alpha- and beta-band bursts significantly shorter than their motor-cortex counterparts ($p < 0.05$, Wilcoxon rank-sum, $N=6$). The method's high specificity makes it practical for closed-loop neuromodulation, neurofeedback, and biomarker discovery in disorders where waveform asymmetry matters.
\end{abstract}

\section{Introduction}

Neural oscillations in the mammalian brain are thought to coordinate activity across distributed regions and to support the integration of sensory input, the control of movement, and the maintenance of cognitive state \citep{buzsaki2004neuronal,fries2015rhythms,giraud2012cortical,jensen2010shaping}. Detecting these oscillations is therefore a foundational step in much of systems and cognitive neuroscience: their presence, duration, location, and fundamental frequency each carry information about the underlying neural process. In practical terms, three questions are routinely asked of a recording. \emph{When} does an oscillation occur, i.e. what are its onset and offset times? \emph{Where} is it generated, i.e. which brain region or layer gives rise to it? And \emph{what} is its fundamental frequency, i.e. which band does it truly occupy \citep{buzsaki2004neuronal,pfurtscheller1999eventrelated}? Answering the \emph{when}, \emph{where}, and \emph{what} is a prerequisite for more ambitious questions about the \emph{why}, the \emph{how}, and the \emph{whom} (the pathology) of neural rhythms.

The dominant analytical tool for answering these questions is the Fast Fourier Transform (FFT) and its derivatives. FFT-based methods decompose a signal into sinusoidal components and then search for peaks over a $1/f$ aperiodic background, a strategy made rigorous by the specparam/FOOOF framework \citep{donoghue2020parameterizing} and extended to time-frequency by SPRiNT \citep{wilson2022sprint} and to discrete oscillation events by OEvent \citep{neymotin2022oevent}. These methods assume, implicitly or explicitly, that the underlying oscillations are sinusoidal. A growing body of work, however, shows that neural oscillations are routinely non-sinusoidal: mu rhythms, hippocampal theta, and cortical beta all exhibit systematic asymmetries between peaks and troughs \citep{belluscio2012crossfrequency,cole2017waveform,cole2019cycle,donoghue2022methodological}. When an FFT is applied to a non-sinusoidal rhythm, the asymmetry of the waveform is expressed as a cascade of harmonic peaks at integer multiples of the fundamental frequency. These harmonic peaks are artifacts of the sinusoidal basis, yet they appear indistinguishable from independent oscillations in the power spectrum. The result is an inflated false-positive detection rate that contaminates downstream analyses and, in some cases, overturns physiological interpretation.

The ambiguity is concrete. An unfiltered electrocorticographic recording from auditory cortex can display peaks over $1/f$ at 6, 12, and 18~Hz. Without further analysis, it is impossible to tell whether these are three independent rhythms or a single 6~Hz non-sinusoidal oscillation with two harmonics. Phase-phase coupling analysis \citep{belluscio2012crossfrequency} can resolve the question: if the 18~Hz component is phase-locked to the 6~Hz component in a 1:3 ratio, it is a harmonic. But the permutation tests required by this approach scale poorly across channel counts, trial counts, and frequency pairs. Phase-phase coupling is therefore impractical for the real-time and closed-loop settings in which oscillation detection is most needed, such as brain-computer interfaces, neurofeedback systems, and phase-locked deep brain stimulation \citep{cagnan2017phase,schalk2004bci2000}.

In this study, we define the principle criteria that characterize a neural oscillation and synthesize them into an algorithm that simultaneously recovers the \emph{when}, the \emph{where}, and the \emph{what} of non-sinusoidal oscillations. The method, which we call Cyclic Homogeneous Oscillation (CHO) detection, combines three ideas. First, following FOOOF, CHO uses a parametric $1/f$ estimate to threshold the time-frequency map and generate candidate bounding boxes---refining the empirical thresholding of OEvent \citep{neymotin2022oevent}. Second, CHO rejects any candidate that exhibits fewer than two full cycles, thereby excluding evoked and event-related potentials whose spectral signatures mimic theta or alpha oscillations \citep{decheveigne2019filters}. Third, CHO tests each surviving bounding box with auto-correlation: the positive-peak spacing of the auto-correlation of the raw signal within the box must agree with the center frequency of the box. Because auto-correlation reads periodicity directly from the time-domain signal, it is insensitive to waveform shape, and the test rejects the harmonic peaks that plague FFT-based approaches. A box whose center frequency is 24~Hz but whose raw signal periodicity is 7~Hz is rejected. We evaluate CHO on synthetic non-sinusoidal bursts (2.5 cycles, 1--3~s long) convolved with $1/f$ noise, and on four human datasets: ECoG from 8 subjects and scalp EEG from 7 subjects during a pre-stimulus auditory reaction-time task, resting-state ECoG (6 subjects), and resting-state SEEG from 6 subjects with depth electrodes in the hippocampus. Across these datasets we compare CHO to FOOOF \citep{donoghue2020parameterizing}, OEvent \citep{neymotin2022oevent}, and SPRiNT \citep{wilson2022sprint}.

\section{Related Work}
\label{sec:relwork}

\paragraph{Frequency-domain detection.} The FFT and its descendants separate a recording into sinusoidal components within canonical bands. Peaks over a $1/f$ aperiodic floor are interpreted as oscillations. FOOOF (now specparam) formalized this by jointly fitting a Lorentzian-like aperiodic component and a mixture of Gaussian peaks \citep{donoghue2020parameterizing}, and it is the de facto standard for peak-frequency estimation. Its principal limitation in our setting is that it operates only on the average spectrum and has no native mechanism for rejecting harmonic peaks.

\paragraph{Time-frequency and event-based detection.} SPRiNT extends FOOOF into the time-frequency domain by parameterizing short-time spectra \citep{wilson2022sprint}. OEvent identifies bounding boxes of elevated power in a wavelet map as discrete oscillation events and summarizes each by center frequency, duration, and number of cycles \citep{neymotin2022oevent}. Neither approach, however, tests whether the bounding box corresponds to a genuine fundamental frequency or to a harmonic of a lower-band rhythm.

\paragraph{Time-domain and waveform-based detection.} A separate strand of work studies oscillations directly in the time domain, quantifying peak-to-trough asymmetry, rise-decay asymmetry, and instantaneous frequency \citep{cole2017waveform,cole2019cycle,quinn2021within,gips2017lineclock}. These methods are well-suited to describing waveform shape \emph{after} an oscillation has been identified, but they do not themselves adjudicate whether a candidate spectral peak is fundamental or harmonic.

\paragraph{Cross-frequency coupling.} \citet{belluscio2012crossfrequency} established that $n{:}m$ phase-phase coupling can identify harmonic relations, and follow-up work has used this tool to disentangle theta-gamma and theta-beta couplings in hippocampus and cortex. The method is statistically principled but computationally heavy, requiring surrogate distributions across all channel-frequency pairs, which puts it out of reach of online analysis pipelines.

\paragraph{The gap.} No existing method jointly (i)~detects the onset and offset of a rhythm, (ii)~rejects harmonic peaks without resorting to permutation-based cross-frequency tests, and (iii)~scales to real-time analysis suitable for closed-loop neuromodulation. CHO is designed to fill this gap by combining bounding-box detection with an auto-correlation periodicity test, at a computational cost that is dominated by a single FFT and auto-correlation per candidate box.

\section{Methods}

\subsection{The CHO detection criteria}

CHO identifies neural oscillations that simultaneously satisfy three criteria. \emph{Criterion 1}: an oscillation must manifest as a peak over the aperiodic $1/f$ noise floor in both the time domain and the frequency domain. Operationally, CHO fits a $1/f$ model to the average power spectrum, following the specparam/FOOOF formulation \citep{donoghue2020parameterizing}, and uses the fit as a parametric threshold to identify candidate bounding boxes in the time-frequency map. This replaces the empirical, data-driven threshold used by OEvent \citep{neymotin2022oevent} and makes the rejection of aperiodic power principled rather than hand-tuned. \emph{Criterion 2}: an oscillation must contain at least two full cycles. Evoked potentials (EPs), event-related potentials (ERPs), and spike transients typically produce fewer than two cycles of fluctuation, yet their large amplitudes can generate apparent oscillatory power in the theta and alpha bands \citep{decheveigne2019filters,donoghue2020parameterizing}. Rejecting sub-two-cycle boxes excludes these confounds. \emph{Criterion 3}: the periodicity revealed by the auto-correlation of the raw signal inside a bounding box must agree with the box's spectral center frequency. Auto-correlation measures the similarity of a signal to a delayed version of itself; the spacing of its positive peaks encodes the fundamental period. Because auto-correlation does not assume a sinusoidal basis, its periodicity is unaffected by waveform asymmetry and therefore returns the fundamental frequency even when the FFT splits the energy across harmonic peaks. A box whose center frequency is 24~Hz but whose auto-correlation periodicity corresponds to 7~Hz fails Criterion 3 and is rejected. Finally, the surviving boxes that neighbor each other in frequency and overlap by more than 75\% in time are merged into single oscillation events, following the convention established by \citet{neymotin2022oevent}.

\subsection{Subjects and ethics}

Four cohorts contributed data. Eight human subjects implanted with subdural ECoG electrodes (x1--x8, 4 females, average age $41 \pm 14$ years, average IQ $96 \pm 18$, range 75--120) participated in an auditory reaction-time task at the Albany Medical Center. Platinum-iridium ECoG electrodes were 4~mm in diameter (2.3~mm exposed) and spaced 10~mm center-to-center, embedded in silicone. Seven healthy subjects (y1--y7, all males, average age $27 \pm 3.6$~years) provided scalp EEG during the same task as a control group. Six subjects contributed resting-state ECoG (ze1--ze8, $N=6$). Six subjects (zs1--zs6, 3 females, average age $46 \pm 16.6$~years) implanted with stereo-EEG depth electrodes participated in resting-state recordings at Barnes Jewish Hospital in St.~Louis; SEEG platinum-iridium electrodes were 2~mm in length (0.8~mm diameter) with 3.5--5~mm center-to-center spacing. For each subject, postoperative CT images were co-registered with preoperative MRI to construct subject-specific cortical models and derive electrode locations \citep{coon2016method}. The Versatile Electrode Localization Framework (VERA) was used to assign SEEG contacts to Brodmann areas during the resting-state analysis.

\subsection{Acquisition and preprocessing}

Raw, unfiltered signals were acquired with the BCI2000 software platform \citep{schalk2004bci2000} interfaced with eight 16-channel g.USBamp devices for EEG, one 256-channel g.HIamp device for ECoG (g.tec., Graz, Austria), and one Nihon Kohden JE-120A long-term recording system for SEEG (Nihon Kohden, Tokyo, Japan). The sampling rate was 1200~Hz for EEG and ECoG and 2000~Hz for SEEG. A clinical safety splitter duplicated ECoG signals to the clinical monitoring system during experimental sessions. Data were then high-pass filtered at 0.05~Hz with a second-order Butterworth filter to remove offsets, and any channel with 60~Hz line noise exceeding ten times the median absolute deviation was excluded from the common-mode reference. To improve the signal-to-noise ratio and reduce computational cost, signals were downsampled via MATLAB's polyphase antialiasing \texttt{resample} to 400~Hz (EEG, ECoG) or 500~Hz (SEEG). For the reaction-time task we analyzed the 1.5~s pre-stimulus period per trial; each ECoG subject performed approximately 400 trials and each EEG subject 200 trials. EEG electrode labels followed the 10--10 system. ECoG channels were assigned to cortical regions using FreeSurfer \citep{fischl2012freesurfer}, and EEG channels were grouped by underlying cortex (e.g., O1 and O2 for visual cortex).

\subsection{Task}

Subjects performed a simple auditory reaction-time task, pressing a button as quickly as possible in response to a 1~kHz pure tone that remained present until the button press. A warning tone penalized responses faster than 100~ms post-stimulus as false starts. Trials proceeded until each subject completed the target trial count. For the motor task analyses, the 1~s interval centered on stimulus onset defined the auditory task window, while the 1~s window spanning 0.5~s before to 0.5~s after the button press defined the motor task window.

\subsection{Phase-phase coupling and time-frequency analysis}

To demonstrate phase-locking between theta- and beta-band components of a single non-sinusoidal oscillation, we used the $n{:}m$ phase-phase coupling estimator of \citet{belluscio2012crossfrequency}. The analysis confirmed the 1:3 phase-locking between 6~Hz theta and 18~Hz beta components in the illustrative auditory cortex example. In evaluating CHO's spectral accuracy, we used four windows yielding a 250~ms estimation period per $1/f$ spectrum. An oscillation detection was counted as accurate when the difference between the ground-truth peak frequency and the detected frequency was less than 1.5~Hz. While CHO currently uses the Hilbert transform for the time-frequency map---because of its MATLAB-native implementation and its control over spectral shape---the pipeline is modular: a wavelet, short-time FFT, Gabor, matching-pursuit, state-space, or empirical mode decomposition front-end \citep{quinn2021within} can be substituted without altering the three-criterion core.

\subsection{Comparison methods}

We compared CHO against three published methods. FOOOF (specparam) \citep{donoghue2020parameterizing} serves as the canonical frequency-domain peak detector. SPRiNT \citep{wilson2022sprint} extends FOOOF into the time-frequency domain by parameterizing short-time spectra. OEvent \citep{neymotin2022oevent} identifies discrete oscillation events as bounding boxes in the time-frequency map. For peak-frequency evaluation we used all three; for onset/offset evaluation the comparison was restricted to OEvent, since FOOOF and SPRiNT do not estimate oscillation boundaries for short oscillatory bursts.

\begin{figure}[t]
  \centering
  \includegraphics[width=0.95\linewidth]{figures/fig_pipeline.pdf}
  \caption{The three-step CHO detection pipeline. Raw electrophysiological signals are first analyzed in the time-frequency domain, where a parametric $1/f$ fit defines bounding boxes around candidate oscillations. Boxes that contain fewer than two complete cycles are rejected to suppress evoked and event-related potentials whose spectral footprint mimics low-frequency oscillations. Surviving boxes are tested against the auto-correlation of the raw signal within each box: the periodicity reported by the positive peaks of the auto-correlation must agree with the center frequency of the box, which rejects harmonic peaks arising from non-sinusoidal waveforms. The output is a set of fundamental oscillations annotated with onset and offset, frequency range, center frequency, and number of cycles.}
  \label{fig:pipeline}
\end{figure}

\begin{figure}[t]
  \centering
  \includegraphics[width=0.95\linewidth]{figures/fig_resting_frequencies.pdf}
  \caption{Predominant fundamental frequency (left) and average oscillation duration (right) detected by CHO across four brain regions during resting state. Primary auditory cortex (Brodmann 41/42) is dominated by 11~Hz oscillations of 450~ms duration; primary motor cortex (Brodmann 4) by 7~Hz oscillations of 450~ms, with additional beta rhythms of $>$500~ms duration; Broca's area by 17~Hz beta oscillations; and hippocampus by 8~Hz oscillations of 450~ms. Values are taken from CHO applied to SEEG and ECoG recordings at rest.}
  \label{fig:resting}
\end{figure}

\begin{figure}[t]
  \centering
  \includegraphics[width=0.85\linewidth]{figures/fig_entropy.pdf}
  \caption{Spatial focality of oscillation detection across 64 EEG channels, quantified by entropy of the per-channel detection distribution. Lower entropy indicates a more focal spatial distribution. The theoretical maximum entropy across 64 channels is $\log_2 64 = 4.16$. In the alpha band, CHO (3.82) produces a substantially more focal distribution than FOOOF (4.15); in the beta band CHO (4.05) is again more focal than FOOOF (4.15). This pattern is consistent with CHO rejecting harmonic peaks that FOOOF detects across spatially diffuse regions.}
  \label{fig:entropy}
\end{figure}

\begin{figure}[t]
  \centering
  \includegraphics[width=0.95\linewidth]{figures/fig_onset_offset.pdf}
  \caption{Onset and offset of 7~Hz oscillations in auditory cortex during an auditory reaction-time task, detected by CHO from ECoG recordings. After the auditory stimulus, subjects press a button on average at 200~ms post-stimulus. CHO-detected offset of 7~Hz oscillations peaks at 300~ms post-stimulus (i.e., 200~ms post-response), and onset re-emergence peaks at 950~ms post-stimulus. Averaged across subjects, 7~Hz oscillations de-synchronize 250~ms and re-synchronize 900~ms post-stimulus.}
  \label{fig:onset_offset}
\end{figure}

\section{Results}

\subsection{Synthetic non-sinusoidal oscillations}

We generated 5~s pink-noise traces (power $\propto 1/f$) and added non-sinusoidal oscillatory bursts of varying length (one cycle, 2.5 cycles, 1~s duration, 3~s duration). Longer bursts produced stronger harmonic peaks, consistent with expectations for non-sinusoidal waveforms. Across the full SNR range, CHO consistently outperformed conventional techniques in specificity and in peak-frequency accuracy. The average specificity of SPRiNT on these signals was below 0.3, reflecting its vulnerability to harmonic peaks and one-cycled transients; CHO's specificity remained robust across the same SNR range. Sensitivity told the reverse story at very low SNR: existing methods were overly sensitive, reporting oscillations where none were present, which is precisely what limited their specificity. At the SNR levels measured in our own recordings ($-7$~dB for EEG and $-6$~dB for ECoG, quantified from Figure~4---figure supplement~1 of the source data), CHO's sensitivity became comparable to SPRiNT's, and its combined accuracy---jointly considering specificity and sensitivity---was the highest of the four methods. The onset/offset evaluation, restricted to OEvent, showed the same pattern: CHO was more specific than OEvent, though slightly less sensitive, in recovering the start and end of synthetic bursts. We also evaluated CHO on purely sinusoidal oscillations, where, as expected, FOOOF and SPRiNT were reasonably specific; CHO retained its specificity advantage even in this easier regime, and its sensitivity scaled with SNR.

\subsection{ECoG during the auditory reaction-time task}

We expected low-frequency oscillations in task-relevant regions---auditory and motor cortex---during the pre-stimulus period. As expected, alpha and beta oscillations appeared in these regions. The topography, however, differed sharply between methods. FOOOF detected alpha and beta oscillations over superior temporal gyrus (STG) and the pre-motor area; CHO detected alpha oscillations over STG and motor cortex but detected beta oscillations exclusively in the pre-motor area. This pattern held across all eight ECoG subjects. Formally, the two methods were statistically comparable in the theta and alpha bands, though CHO reported smaller median occurrence rates than FOOOF. In the beta band, CHO's detection rate was significantly lower than FOOOF's in most regions, with the exception of precentral, pars opercularis, and caudal middle frontal areas (Wilcoxon rank-sum test, $p < 0.05$ after Bonferroni correction). FOOOF also showed higher sensitivity for delta, theta, and low-gamma oscillations, though both methods detected comparatively few events in those bands. The most parsimonious interpretation of the beta discrepancy is that FOOOF's beta detections in regions such as temporal cortex represent harmonics of the predominant theta and alpha rhythms---an interpretation supported by our synthetic results, where harmonic peaks inflate FOOOF's false-positive rate.

\subsection{EEG during the auditory reaction-time task}

The EEG results recapitulated the ECoG findings in broad outline. FOOOF identified alpha oscillations across mid-frontal and visual areas and beta oscillations across the whole scalp; CHO localized alpha to the visual cortex and beta to the pre-motor region, producing a markedly more focal topography (Figure~\ref{fig:entropy}). To quantify focality, we computed the entropy of the detection distribution across 64 channels; lower entropy indicates a more concentrated spatial pattern, with a theoretical maximum of $\log_2 64 = 4.16$. In the alpha band, CHO's entropy was 3.82 versus FOOOF's 4.15; in the beta band, 4.05 versus 4.15. CHO detected more alpha oscillations at the occipital electrodes (O1, O2) than at the pre-motor electrodes (FC2, FC4), whereas FOOOF reported alpha and beta more evenly across the scalp, including higher beta at FC2 and FC4 than at O1/O2. Consistent with ECoG, FOOOF was more sensitive in delta, theta, and low-gamma bands; unlike ECoG, low-gamma detections were more common in EEG, which we attribute to the greater spectral mixing of scalp signals.

\subsection{Onset and offset of oscillations}

A central motivation for CHO is the recovery of oscillation boundaries (the \emph{when}). Applied to ECoG channels with a significant broadband gamma response to the auditory stimulus, CHO detected 7~Hz oscillations whose onset and offset distributions were uniform during the 1.5~s pre-stimulus period---consistent with stochastic, task-unrelated bursting---and became Gaussian after the stimulus (Figure~\ref{fig:onset_offset}). Subjects pressed the button 200~ms post-stimulus on average. The peak in detected offsets occurred at 300~ms post-stimulus (200~ms post-response), and the peak in detected onsets at 950~ms post-stimulus. Averaged across subjects, the 7~Hz oscillations de-synchronized at 250~ms and re-synchronized at 900~ms post-stimulus. The 7~Hz oscillations are thus suppressed during stimulus processing and response, and re-emerge only after the motor act is complete. This temporal profile is consistent with event-related desynchronization \citep{pfurtscheller1999eventrelated} and provides a concrete illustration of CHO's ability to resolve the \emph{when} in addition to the \emph{where} and \emph{what}.

\subsection{SEEG hippocampal oscillations and harmonics}

Hippocampal oscillations are known to be non-sinusoidal and of particular interest to memory research \citep{belluscio2012crossfrequency}. We recorded from SEEG electrodes placed in the left anterior hippocampus during resting state. The power spectrum showed neural activity over the $1/f$ trend up to 30~Hz, but offered little specificity about the underlying rhythms. CHO, in contrast, identified distinct fast theta bursts (see Figure~\ref{fig:pipeline} for the staged processing; the raw power spectrum indicated broadband activity whose auto-correlation revealed a cleaner fundamental structure). Consistent with the harmonic interpretation, the non-sinusoidal alpha rhythm produced beta-band harmonics in the time-frequency map that CHO's Criterion 3 correctly rejected. CHO found primarily alpha-band oscillations in hippocampus, with beta-band peaks clustered near the alpha band (8--14~Hz) rather than in the conventional 13--30~Hz window. Cross-subject comparison of CHO and FOOOF in hippocampus did not reach statistical significance due to the small sample ($N=6$) and variable electrode placement, but CHO systematically produced less overlap between alpha and beta bands and essentially no low-gamma (30--40~Hz) detections.

\subsection{Resting-state frequency and duration across regions}

Using VERA to assign SEEG contacts to Brodmann areas, we characterized the dominant oscillation frequency and duration in four regions (Figure~\ref{fig:resting}). In primary auditory cortex (Brodmann areas 41 and 42), CHO found 7 and 11~Hz as the predominant frequencies, with the 11~Hz rhythm lasting 450~ms on average. In primary motor cortex (Brodmann area 4), the predominant fundamental frequency was 7~Hz at 450~ms duration, and the motor cortex showed more beta-band oscillations ($>$500~ms duration) than auditory cortex. Broca's area mirrored motor cortex in profile but with a predominant beta frequency of 17~Hz, lower than the 22 or 24~Hz beta rhythms typical of motor cortex. Hippocampus, finally, was dominated by 8~Hz oscillations of 450~ms duration. The alpha- and beta-band bursts in hippocampus were significantly shorter than those in motor cortex at rest ($p < 0.05$, Wilcoxon rank-sum, $N=6$), consistent with the view that hippocampal rhythms occur as brief, repeated fast-theta bursts rather than sustained oscillations.

\section{Discussion}

\paragraph{Why three criteria.}
CHO's high specificity is a direct consequence of the three criteria it enforces. The first criterion---a peak over the parametric $1/f$ floor---prevents aperiodic power from being misread as oscillatory activity, a confound that \citet{donoghue2020parameterizing} showed to contaminate classical band-limited analyses. The second criterion---at least two full cycles---removes evoked and event-related potentials, whose broadband spectral signatures can closely resemble theta or alpha peaks. The third criterion---agreement between spectral center frequency and the periodicity of the auto-correlation of the raw signal---rejects harmonic peaks. It is the third criterion that most clearly distinguishes CHO from prior approaches. FOOOF, SPRiNT, and OEvent all read frequency from a spectral representation and therefore accept harmonic peaks as independent rhythms; CHO demands that the time-domain periodicity match, which is a test of waveform structure rather than of spectral energy. The criterion is cheap to evaluate---one auto-correlation per candidate box---which matters for closed-loop applications where a single phase estimate must be produced within tens of milliseconds \citep{cagnan2017phase,schalk2004bci2000}.

\paragraph{Focal beta and the harmonic interpretation.}
The most striking difference between CHO and FOOOF in our empirical data concerns beta oscillations. FOOOF reports beta broadly across the scalp and cortex; CHO reports beta almost exclusively in pre-motor and frontal areas prior to the button press (Figure~\ref{fig:entropy}). These regions are long-standing locations of movement-related beta activity \citep{pfurtscheller1999eventrelated,jenkinson2011new}. The beta detections that FOOOF produces elsewhere---in temporal cortex, in occipital cortex---are, on our account, harmonics of the local alpha rhythm, not independent oscillations. The same logic applies to low-gamma, which FOOOF detects in both ECoG and EEG in regions where CHO does not; if the low-gamma peaks are harmonics of a non-sinusoidal beta, CHO's Criterion 3 will correctly suppress them while FOOOF has no grounds to. We are not arguing that beta and gamma rhythms never occur outside premotor cortex; the claim is narrower: when the local waveform is sharply non-sinusoidal, the FFT-based detections outside the canonical generators are likely to be harmonics. Distinguishing the two cases requires a method whose acceptance criterion is insensitive to waveform shape, which is exactly what auto-correlation provides.

\paragraph{Hippocampal rhythms and the wide-band view.}
The hippocampal results support a growing consensus that hippocampal oscillations cannot be cleanly partitioned into canonical bands. Recent studies report hippocampal rhythms up to 40~Hz \citep{cole2019cycle} or in the 3--12~Hz range \citep{cole2017waveform}, and the mismatch between these reports and the textbook theta/alpha ranges \citep{buzsaki2004neuronal} reflects the intrinsic non-sinusoidality of the underlying activity. CHO's wide-band analysis with auto-correlation-based periodicity testing dissolves this tension: the apparent alpha-beta split in hippocampus is, on our data, a single non-sinusoidal rhythm whose harmonics land in the beta band. More generally, the same phenomenon can occur in any region where the fundamental oscillation is asymmetric, and our ECoG and EEG analyses are consistent with this. The practical implication is that oscillation analyses, especially in hippocampus and motor cortex, should be performed on wide-band signals and with waveform-aware tools \citep{donoghue2022methodological,quinn2021within}, not on band-pass-filtered traces within canonical bins.

\paragraph{Applications: closed-loop neuromodulation and neurofeedback.}
Knowing when an oscillation begins, how long it lasts, and what its fundamental frequency is, is a practical as well as a scientific question. Phase-locked deep brain stimulation---delivering stimulation pulses at a specific phase of an ongoing oscillation---has been shown to reduce the total stimulation required to achieve a therapeutic effect \citep{cagnan2017phase}, with consequences for battery life, surgical burden, and side-effect profile. Realizing phase-locked stimulation in practice demands high-specificity detection: a false-positive oscillation triggers a useless pulse, and a false-positive frequency delivers the pulse at the wrong phase. CHO's combination of bounding-box onset/offset detection with harmonic-rejecting periodicity testing provides exactly this profile. Similar considerations apply to neurofeedback, where the targeted rhythm must be distinguished reliably from unrelated activity if the feedback signal is to be meaningful \citep{jensen2010shaping}. The modular design of CHO means the time-frequency front-end can be swapped without touching the three-criterion core, which makes it straightforward to adapt the method to low-latency hardware.

\paragraph{Limitations and routes to higher sensitivity.}
CHO favors specificity over sensitivity at low SNR. At our physiological SNR levels ($-7$~dB for EEG, $-6$~dB for ECoG), CHO recovers 50--60\% of true oscillations. This is adequate for mapping and for closed-loop control, where false positives are the more expensive error, but it is not ideal for detection tasks where every oscillation matters. Several extensions could raise sensitivity without sacrificing specificity. A wavelet or Gabor time-frequency front-end could improve the localization of short bursts. Replacing the FFT-based periodicity estimate with a state-space model or empirical mode decomposition \citep{quinn2021within} could stabilize the auto-correlation in the presence of strong aperiodic activity. Averaging across spatially correlated contacts---for instance, neighboring SEEG sites within hippocampus---could improve SNR directly. A separate limitation is that CHO characterizes the fundamental frequency of a non-sinusoidal oscillation but does not yet describe its harmonic structure. Harmonic structure carries information about pathology \citep{cole2017waveform}, and future work that combines CHO's fundamental-frequency estimate with a harmonic-profile estimate could yield biomarkers for Parkinson's disease, epileptic seizure-onset zones, and other conditions where waveform asymmetry is diagnostic.

\section{Conclusion}

Neural oscillations coordinate activity across brain regions, and detecting them well is foundational to systems, cognitive, and clinical neuroscience. We presented Cyclic Homogeneous Oscillation (CHO) detection, a method built on three criteria---a peak over the parametric $1/f$ floor, at least two cycles, and agreement between spectral center frequency and auto-correlation periodicity---that together reject the harmonic peaks and sub-threshold transients that contaminate conventional detectors. On synthetic non-sinusoidal bursts at physiological SNR, CHO is more specific than FOOOF, OEvent, and SPRiNT while remaining sensitive enough to recover 50--60\% of oscillations. On human ECoG, EEG, and SEEG recordings, CHO produces more focal topographies than FOOOF (alpha entropy 3.82 vs.~4.15; beta entropy 4.05 vs.~4.15 across 64 EEG channels), resolves the onset/offset of 7~Hz auditory-cortex oscillations around a reaction-time task, and recovers regionally characteristic fundamental frequencies at rest: 11~Hz in auditory cortex, 7~Hz in motor cortex, 17~Hz in Broca's area, 8~Hz in hippocampus, each of duration approximately 450~ms. The method is modular, computationally light, and directly compatible with closed-loop neuromodulation and neurofeedback. Its high specificity makes it a practical substrate for the identification of oscillatory biomarkers in disorders where non-sinusoidal waveform structure is diagnostic.

\bibliography{references}

\end{document}
